% Transcription — fonds Alexandre Grothendieck, shelfmark 19, batch 4 (pages 61-80).
% Established from the facsimile archives/batches/19/batch-04.pdf (a mirror of
% https://grothendieck.umontpellier.fr/19.pdf), per the skill
% transcribe-grothendieck.
%
% Page 69 is a blank sheet — absent.
%
% Three pieces. Pages 61-68 close the "Théorèmes de descente de M. Artin –
% P. Cartier" begun at page 54 of batch 3, and exhaust the "11 p" the page 54
% wrapper counts (pages 56 to 66). Grothendieck paginates them himself: page 61
% is his 5, page 62 his "5 bis", page 64 his "5 ter", page 66 his 6. From page
% 62 to page 67 he writes in English, and the transcription keeps his language;
% only the apparatus is French. Pages 62, 63, 65 and 66 are cancelled by long
% diagonal strokes over the whole sheet. Page 67 carries a census of the twelve
% functors and the two exact sequences that go with the open/closed picture;
% page 68, back in French, is a note on the limits the pair (i*, j*) commutes
% with.
%
% Page 70 is a piece on its own — one sheet numbered 1 in a circle — deriving
% the co-semi-simplicial object Phi*(X) from an adjunction f, g, with phi = fg,
% lambda = f * alpha * g, and the coaugmentation shown to be a homotopy
% equivalence.
%
% Page 71 is the wrapper of the third piece and names it: "Fidélité". Pages
% 72-75 are the theory of transportable faithful functors read as forgetful
% functors: p : E -> B transportable, p fidèle iff the fibres are ordered
% categories, p conservatif iff they are groupoids, the four equivalent
% conditions of page 74 (fibrant + fidèle + conservatif, fibres discrètes,
% E_{/X} -> B_{/p(X)} an equivalence, E ≈ B_{/F}), and the reconstruction of E
% from the functor Sigma : B° -> (ens. ordonnés). Page 73 carries a single
% pencil line.
%
% Page 76 is the wrapper of the fourth piece — "Morphismes et comor[phismes] —
% Liste Problèmes". Pages 77-80 are the list itself, "Questions sur les topos",
% eleven numbered questions on immersions of topoi, dominant morphisms, fibred
% categories of topoi, and the tower of successive adjoints f_!, f*, f_*, f^!,
% which the page 80 proposition and question take up. Page 80 breaks off
% mid-sentence and runs into batch 5. These four sheets are pencilled and
% heavily overwritten; they are the hardest of the batch and are transcribed
% with many gaps.
%
% The dating is the inventory's own, for the whole shelfmark; nothing on these
% pages carries a date.
%
% Pass: Opus 5 (claude-opus-5), 2026-08-22 — first pass, unchecked against the pages by a human.
\documentclass[11pt,a4paper]{article}
% Le préambule commun à toutes les transcriptions du fonds Grothendieck.
%
% Il tient en une idée : **l'incertitude se compose, elle ne se résout pas.**
% Une transcription qui lisse un mot douteux a détruit la seule chose pour
% laquelle on la faisait — un lecteur doit pouvoir savoir, à chaque endroit,
% ce qui était sur le feuillet et ce qui a été ajouté. Les six macros qui
% suivent sont donc l'appareil critique, et non de la décoration.
%
% Les mêmes commandes sont reconnues par `scripts/render.mjs`, qui produit la
% vue de lecture du site. Une macro ajoutée ici doit l'être aussi là-bas,
% faute de quoi le rendu échouera bruyamment — ce qui est voulu : un rendu
% silencieusement faux est pire qu'un rendu absent.

\ProvidesPackage{grothendieck}[2026/08/08 Transcriptions du fonds Grothendieck]

\usepackage{amsmath,amssymb,amsthm}
\usepackage{mathrsfs}
\usepackage{stmaryrd}
% Les diagrammes commutatifs sont partout dans ce fonds : c'est la langue
% même de la géométrie algébrique de Grothendieck, et une transcription qui
% les rendrait en prose ne transcrirait rien.
\usepackage{tikz-cd}
% Français par défaut — c'est la langue des feuillets — mais l'anglais est
% chargé aussi : la traduction et le résumé sont des documents anglais, et la
% typographie française (espaces avant les deux-points) y serait une faute.
% Ces documents basculent par \selectlanguage{english} en tête de corps.
\usepackage[english,french]{babel}
\usepackage{fontspec}
\usepackage{geometry}
\geometry{a4paper,margin=25mm,marginparwidth=28mm}
\usepackage{marginnote}
\usepackage{xcolor}
% ulem, not soul. `soul`'s \st and \ul are built on a character-by-character
% scan that XeLaTeX's font machinery breaks: the document fails with an
% undefined control sequence rather than degrading. ulem does the same two jobs
% and is Unicode-engine safe. `normalem` stops it hijacking \emph, which the
% transcriptions use for Grothendieck's own emphasis.
\usepackage[normalem]{ulem}
\usepackage{hyperref}
\hypersetup{colorlinks=true,linkcolor=black,urlcolor=black,pdfborder={0 0 0}}

% --- Métadonnées du lot -----------------------------------------------------
% Reprises telles quelles par le script de rendu, qui les affiche en tête de
% la vue de lecture. Elles fixent ce que le fichier prétend couvrir : sans
% elles, une transcription flotte sans point d'attache dans un fonds de seize
% mille feuillets.

\newcommand{\folder}[1]{\gdef\grothFolder{#1}}
\newcommand{\batch}[1]{\gdef\grothBatch{#1}}
\newcommand{\leaves}[2]{\gdef\grothFirst{#1}\gdef\grothLast{#2}}
\newcommand{\foldertitle}[1]{\gdef\grothFoldertitle{#1}}
\gdef\grothFolder{?}\gdef\grothBatch{?}\gdef\grothFirst{?}\gdef\grothLast{?}\gdef\grothFoldertitle{}

% --- Le résumé --------------------------------------------------------------
%
% Il ouvre la lecture modernisée, sous le titre « Résumé », et il est composé
% en petit. Ce n'est pas un ornement : c'est le seul passage du document qui
% ne soit pas des mathématiques, et un lecteur doit pouvoir le reconnaître
% avant de l'avoir lu — puis le sauter s'il connaît déjà le sujet, ou s'y
% arrêter s'il ne le connaît pas. Le filet qui le suit marque la reprise du
% corps.

\newenvironment{resume}
  {\small\setlength{\parskip}{0.45em}}
  {\par\medskip\hrule\medskip}

% --- Mots-clés ---------------------------------------------------------------
%
% En anglais, en clôture du résumé de la lecture modernisée : le vocabulaire
% moderne sous lequel chercher ce que les feuillets construisent. Le manifeste
% du site (`npm run manifest`) les extrait pour étiqueter la cote entière —
% cette ligne est la source unique des tags, il n'y a pas de fichier à côté.
\newcommand{\keywords}[1]{\par\smallskip\noindent
  {\footnotesize\textsc{Keywords} --- \textit{#1}}\par}

% --- L'appareil critique ----------------------------------------------------

% Le numéro de feuillet, en marge. C'est l'ancre qui permet de régler un
% désaccord sur une formule en regardant *un* feuillet plutôt qu'une chemise
% de six cents. Le site s'en sert aussi pour tourner les pages du fac-similé
% quand on fait défiler la transcription.
\newcommand{\leaf}[1]{%
  \marginnote{\footnotesize\color{gray!70}\textsf{#1}}%
  \label{leaf:#1}\ignorespaces}

% Illisible : on ne devine pas. Un mot inventé qui se lit comme les autres est
% le pire résultat possible.
\newcommand{\ill}{\textcolor{red!60!black}{[\,\ldots\,]}}

% Lecture douteuse : le mot est donné, et signalé comme incertain.
\newcommand{\uncertain}[1]{\uline{#1}}

% Ajout de l'éditeur — une lettre manquante, un mot sous-entendu.
\newcommand{\add}[1]{\textcolor{blue!50!black}{[#1]}}

% Biffé par Grothendieck lui-même. Ce qu'il a rayé fait partie du document :
% une rature dit souvent où la pensée a changé de direction.
\newcommand{\struck}[1]{\sout{#1}}

% Note du transcripteur, sur l'état matériel du feuillet.
\newcommand{\note}[1]{\par\smallskip{\small\color{gray!60!black}\textit{[#1]}}\par\smallskip}

% Note marginale de Grothendieck, distincte du corps du texte.
\newcommand{\marginal}[1]{\par\smallskip\noindent\hspace{1em}%
  {\small\color{gray!60!black}#1}\par\smallskip}

% --- Titre ------------------------------------------------------------------

\AtBeginDocument{%
  \begin{center}
    {\large\bfseries Fonds Alexandre Grothendieck --- cote n\textsuperscript{o}~\grothFolder}\\[2pt]
    {\itshape\grothFoldertitle}\\[4pt]
    {\small feuillets \grothFirst--\grothLast\ (lot \grothBatch)}\\[2pt]
    {\footnotesize\color{gray!60!black}Transcription établie d'après le fac-similé numérisé
     par l'Université de Montpellier.\\
     Les lectures incertaines sont soulignées, les illisibles notées [\,\ldots\,],
     les ajouts entre crochets.}
  \end{center}
  \vspace{1em}}

\endinput


\folder{19}
\batch{4}
\pages{61}{80}
\dating{[à partir de 1958-1973]}
\watermark{Édition de démonstration}
\foldertitle{Topos : notes manuscrites (s.d.)}

\begin{document}

\section*{Théorèmes de descente de M. Artin – P. Cartier (suite, pages 61 à 68)}

\page{61}\note{l'auteur numérote cette page 5, dans la suite commencée page 56}

\textbf{Cor.\ 2} Sous les conditions de la prop.\ 1, conditions équivalentes :
\note{une phrase insérée entre les lignes, « néc.\ \ill{} », n'a pas pu être
lue}

\begin{enumerate}
\item[a)] $j^{*}$ existe \add{et est exact}, essentiellement \note{« néc.\ pl.\ fid.\ » est inséré au-dessous} ;
\item[b)] $i_{!}$ existe, et a une image épaisse ;
\item[b$'$)] $i_{!}$ existe, et $i_{!} i^{*} F \to F$ est injectif pour tout $F$.
\end{enumerate}

Dans ce cas, $i_{!} = \mathrm{Ker}\, j^{*}$, $j^{*} = \mathrm{Coker}\, i_{!}$,
\ill{} \uncertain{foncteur} $i_{!}$ exact. [N.B.\ Mais il ne \emph{suffit} pas
que $j^{*}$, $i_{!}$ existent, $i_{!}$ exact \ill{}]

La prop.\ 1 et le corollaire prouvent que \ill{} \ill{} pour b), b$'$).

\page{62}\note{l'auteur numérote cette page « 5 bis » ; toute la page est
barrée de longs traits diagonaux}

$C \longrightarrow D$

\begin{tikzcd}[column sep=huge]
C'' \arrow[r, bend left=25, "j_{*}"] &
C \arrow[l, bend left=25, "j^{*}"] \arrow[r, "i^{*}"] &
C'
\end{tikzcd}

\begin{enumerate}
\item[a)] $i_{!}$ exists, and has thick image [i.e.\ $i_{!} i^{*} F \to F$
      injective] ;
\item[b)] $j^{*}$ exists and is exact.
      \marginal{[N.B.\ even if $i_{!}$ exists and is exact, so that $j^{*}$
      exists, we don't see that $j^{*}$ \ill{} be exact]}
\end{enumerate}

\struck{If so, $i_{!} = \mathrm{Ker}\, j^{*}$, $j^{*} \simeq \mathrm{Coker}\,
i_{!}$.}

We have seen already a) $\Longrightarrow$ b), let us prove b)
$\Longrightarrow$ a). Therefore, let us define the functor
\[
  \tau(F) = \mathrm{Ker}\bigl(F \longrightarrow j_{*} j^{*} F\bigr) ,
\]
we see that $\tau$ is an \struck{\ill{}} exact functor [thanks to the fact that
$j_{*} j^{*} F$ is exact] zero on \ill{} of the type $j_{*}F$,
\struck{\ill{}} hence we have
\[
  \tau(F) \simeq i_{!} i^{*} F
\]
with a well-defined \struck{exact} functor $i_{!} : C'' \to C$.\note{sic : sur
cette page comme sur les suivantes $i_{!}$ va de $C'$ dans $C$} More \ill{}
\marginal{\ill{}}
\[
  \mathrm{Hom}(i_{!}G, F) \simeq \mathrm{Hom}(G, i^{*}F) ,
\]
and we \uncertain{need} to obtain only
\[
  \mathrm{Hom}(i_{!} i^{*}F', F) \simeq \mathrm{Hom}(i^{*}F', i^{*}F)
\]
i.e.\ $\mathrm{Hom}(\tau F', F) \simeq \mathrm{Hom}(i^{*}F', i^{*}F)$.

But as $j_{*} j^{*} F$ is the biggest \struck{\ill{}} quotient of $F$
\struck{\ill{}}, $\mathrm{Im}\, j_{*} = \mathrm{Ker}\, i^{*}$, this comes
\ill{} immediately, by the construction of the quotient category $C/C''$.

\textbf{Question.} \struck{Does} The existence of $j^{*}$ (exact or not) does
\struck{\ill{}} imply the existence of $i_{!}$ ?

\page{63}\note{page barrée de longs traits diagonaux}

\[
  \begin{array}{lll}
    j^{*} i_{!} = 0 & i_{!} = \mathrm{Ker}\, j^{*} & j^{*} = \mathrm{Coker}\, i_{!} \\
    i^{*} j_{*} = 0 & j_{*} = \mathrm{Ker}\, i^{*} & i^{*} = \mathrm{Coker}\, j_{*}
  \end{array}
\]
\[
  \begin{array}{ccccccc}
    0 & \to & N & \to & i_{!} i^{*} F & \to & F \\
    0 & \to & i^{*}N & \to & i^{*} i_{!} i^{*} F & \to & i^{*}F \\
    0 & \to & i_{!} i^{*} N & \to & i_{!} i^{*} i_{!} i^{*} F & \to & i_{!} i^{*} F
  \end{array}
\]

\begin{tikzcd}[column sep=huge]
C'' \arrow[r, "j_{*} = j_{!}" description] &
C \arrow[l, bend right=30, "j^{*}"] \arrow[l, bend left=30, "j^{!}"']
  \arrow[r, "i^{*} = i^{!}" description] &
C' \arrow[l, bend right=30, "i_{!}"'] \arrow[l, bend left=30, "i_{*}"]
\end{tikzcd}

2\textsuperscript{e} généralisation : conditions on pairs $i^{*}, i_{!}$

\begin{enumerate}
\item[a)] $i_{!}$ is fully faithful, thick image, and $i_{!} i^{*} F \to F$
      injective ;
\item[b)] $i^{*}$ \uncertain{passage to the quotient} that is exact, and one
      of the following conditions hold (N.B.\ \ill{} by the previous
      assumptions, $j_{*}$ is defined \ill{} said assumptions, and $j^{*}$
      \ill{} must exist \ill{} we checked, and is exact previously \ill{})
      \begin{enumerate}
      \item[$\alpha$)] $i_{!}$ \struck{\ill{}} has thick image ;
      \item[$\beta$)] $i_{!} i^{*} F \to F$ injective for every $F$ ;
      \item[$\delta$)] $i_{!} = \mathrm{Ker}\, j^{*}$.
      \end{enumerate}
\end{enumerate}

[N.B.\ \uncertain{Many} \struck{\ill{}} conditions which are formally
distinct)

This situation is essentially contained in \uncertain{the} exact sequence
\[
  0 \longrightarrow i_{!} i^{*} F \longrightarrow F \longrightarrow
  j_{*} j^{*} F \longrightarrow 0 .
\]

Under these conditions, is it true that these conditions \ill{} are equivalent
to the existence of $j_{!}$ \ill{} ? If $i_{*}$ is \ill{} \struck{\ill{}} it
says
\[
  j_{*} j^{!} F = \mathrm{Ker}\bigl(F \longrightarrow i_{*} i^{*} F\bigr) .
\]
\ill{}

\page{64}\note{l'auteur numérote cette page « 5 ter »}

\marginal{Below}

[Assume moreover that $i^{*}$ exists equally. Then \ill{} situation becomes
interesting. Moreover, in this case, \struck{\ill{}} equivalent conditions
\begin{enumerate}
\item[a)] $j^{!}$ exact $\Longrightarrow$
\item[b)] $i_{*}$ exact $\Longrightarrow$
\item[c)] $\varphi = j^{*} i_{*}$ is exact.
\end{enumerate}
avec, sous b), $\Downarrow$ \uncertain{$j_{*}$} \uncertain{Ker} épaisse
$\Downarrow$ $F \to i_{*} i^{*} F$ surjectif.

\page{65}\note{page barrée de longs traits diagonaux ; le raisonnement est
récrit trois fois par-dessus lui-même et n'est transcrit qu'en partie}

Question: is it enough that $j^{*}$ \ill{} or that $i_{!}$
\struck{be exact} $i_{!} i^{*}F \to F$ have a kernel $N$ such \add{that}
$i^{*}N = 0$ \ill{} Then $j^{*}$ being exact, we have
\[
  0 \to N \to i_{!} i^{*} F \to F , \qquad j^{*}N \to j^{*} i_{!} i^{*} F
  \to j^{*}F .
\]
Yet $j^{*} i_{!} = 0$, hence \ill{} $j^{*}N = 0$ then $N = 0$.
\struck{$i_{!} i^{*} N = N$}

Question. Is it enough $i_{!}$ exact ?? No, as we see in the case dual to :

\begin{tikzcd}[column sep=huge, nodes={font=\scriptsize}]
C'' \arrow[r, bend left=25, "j^{*}"] &
C \arrow[l, bend left=25, "j_{!}"] \arrow[r, bend left=25, "i^{*}"] &
C' \arrow[l, bend left=25, "i_{*}"]
\end{tikzcd}

où $C'' = $ $A$-modules \ill{} killed by some $t \in S$, $C = C_{A}$,
$C' = C_{S^{-1}A}$, $i^{*}$ la localisation et $i_{*}$ la restriction des
scalaires. $i_{*}$ is exact, but not $j_{!}$ in general.
[$A$ \struck{\ill{}} discrète \ill{}, $S = \{\pi\}$, $\pi$ uniformisante)

\page{66}\note{l'auteur numérote cette page 6 ; elle est barrée de longs
traits diagonaux}

\begin{itemize}
\item foncteur (exact, surjectif et $=$ un passage au quotient) $i^{*}$ qui a
      \struck{deux} adjoints $i_{!}$, $i_{*}$ \add{dont l'un} exact, et $=$
      pleinement fid.\ \ill{} (image épaisse) ;
\item foncteur (exact, pleinement fid.\ et $=$ \ill{} image épaisse) $j_{*}$
      qui a deux adjoints \add{$j^{*}$, $j^{!}$} dont l'un \add{$j^{*}$} est
      exact, \struck{surjectif}, et $=$ un passage au quotient) ;
\item situations mixtes \uncertain{$i^{*}, i_{*}, j^{*}, j_{*}$} et autres
      éventuellement ;
\item foncteurs exacts à gauche.
\end{itemize}

Start with situation

\begin{tikzcd}[column sep=huge]
C \arrow[r, bend left=25, "i^{*}"] & C' \arrow[l, bend left=25, "i_{*}"]
\end{tikzcd}

\struck{exact}, $i_{*}$ fully faithful \struck{\ill{}}, $i^{*}$ \struck{exact}
passage to quotient \struck{\ill{}}. Then the following conditions are
equivalent

\begin{enumerate}
\item[a)] $i^{*}$ \struck{exact} has another adjoint $i_{!}$, $i_{!}$
      \struck{is exact} has thick image [N.B.\ $i_{!} i^{*}F \to F$ \ill{} and
      $i_{*}$ \ill{} $i_{!}$ pl.\ fid.] ;
\item[b)] $i^{*}$ is \struck{\ill{}} exact [then of the type $C \to C/j_{*}C'$,
      where $j_{*}$ is the imbedding of a thick subcategory $C''$ of $C$],
      $j^{*}$ exist and is exact ;
\item[c)] equivalent with the "standard" situation defined by \add{an exact}
      functor $\varphi$.
\end{enumerate}

\marginal{Symmetric point of view : start with $C'' \to C$, $j_{*}$ fully
faithful \ill{}, $j^{*}$ passage au quotient $C' = C/C''$ \ill{} a) $j^{*}$
\ill{} adjoint $i^{*} : C \to C'$ \ill{}}

We have already \struck{checked} c) $\Longrightarrow$ a), we have checked
already c) $\Longrightarrow$ b) and b) $\Longrightarrow$ c), now let us check
a) $\Longrightarrow$ b). Of course if we \uncertain{assume} b), then $i^{*}$
is exact \add{having two} adjoints. \struck{\ill{}} I \ill{} $C'' =
\mathrm{Im}$ \ill{}. We want to define $j^{*}F$ :
\[
  j^{*}F = \mathrm{Coker}\bigl(i_{!} i^{*} F \longrightarrow F\bigr) .
\]

\page{67}

\begin{tikzcd}[column sep=huge]
C'' \arrow[r, bend left=25, "j_{*}"] &
C \arrow[l, bend left=25, "j^{*}"] \arrow[r, bend left=25, "i^{*}"] &
C' \arrow[l, bend left=25, "i_{!}"]
\end{tikzcd}

\[
  0 \to F_{U} \to F \to F^{Y} \to 0 , \qquad f_{!} f^{*} , \qquad f = i, j
\]
\[
  0 \to \Gamma_{Y} F \to F \to \Gamma_{U} F , \qquad f_{*} f^{!} , \qquad
  f = i, j
\]

altogether 12 \uncertain{foncteurs} \ill{} :

\begin{enumerate}
\item[] $j_{*}$, $j^{*}$, $j^{!}$ \quad ; \quad $i^{*}$, $i_{!}$, $i_{*}$ ;
\item[] $i_{!} i^{*}$, $i_{*} i^{*}$ \quad $C \to C$ ;
\item[] $j_{*} j^{*}$, $j_{*} j^{!}$ \quad $C \to C$ ;
\item[] $\varphi = j^{*} i_{*}$ \quad $C' \to C''$ ;
\item[] $j_{*} \varphi$ \quad $C' \to C$.
\end{enumerate}

\page{68}

Notons $(i^{*}, j^{*})$ commutent aux $\varprojlim$ (quelconques), et : pour
que $\varprojlim_{i} (G_{i}, H_{i}, \lambda_{i})$ existe, il suffit que
$\varprojlim_{i} G_{i}$ et $\varprojlim_{i} H_{i}$ existent, et alors \ill{}
\marginal{« quelconques » porte au-dessus l'insertion : si elles existent dans
$C'$, $C''$}

Pour les propriétés d'existence \struck{à gauche}, $(i^{*}, j^{*})$ est
\struck{extrêmement} « aussi bon que » $\varphi : C' \to C''$, p.ex.\ si
$\varphi$ commute aux $\varprojlim$ finies (resp.\ quelconques) et si \ill{}
\struck{\ill{}} la $\varprojlim$ existe dans $C'$, $C''$, alors il existe dans
$C$ et les foncteurs $(i^{*}, j^{*})$ y commutent \ldots

\section*{Adjonction et objets co-semi-simpliciaux (page 70)}

\page{70}\note{feuille numérotée 1 dans un cercle}

\begin{tikzcd}[column sep=huge]
A \arrow[r, bend left=25, "f"] & B \arrow[l, bend left=25, "g"]
\end{tikzcd}

$g$ adjoint à \uncertain{gauche} de $f$, donc on a \note{« gauche » est récrit
par-dessus un autre mot}
\[
  \mathrm{id}_{A} \overset{\alpha}{\longrightarrow} gf , \qquad
  fg \overset{\beta}{\longrightarrow} \mathrm{id}_{B} .
\]
On pose $\varphi = fg$. On a aussi
\[
  \varphi \overset{\lambda}{\longrightarrow} \varphi^{2} , \quad\text{i.e.}\quad
  fg \longrightarrow f\,\underbrace{gf}_{\alpha}\,g , \qquad
  \lambda = f * \alpha * g ,
\]
en plus de $\varphi \overset{\beta}{\longrightarrow} \mathrm{id}_{B}$.

$\lambda$ et $\beta$ permettent de définir, pour tout objet $X$ de $B$, une
structure \struck{cosimpliciale} co-semi-simpliciale $\Phi^{*}(X)$ :

\begin{tikzcd}[column sep=large, row sep=small, nodes={font=\scriptsize}]
\varphi(X) \arrow[r, bend left=18, "\varphi * \beta"]
           \arrow[r, bend right=18, "\beta * \varphi"'] \arrow[d, no head] &
\varphi^{2}(X) \arrow[r, bend left=25] \arrow[r] \arrow[r, bend right=25]
           \arrow[d, no head] &
\varphi^{3}(X) \arrow[d, no head] \\
\Phi^{0}(X) & \Phi^{1}(X) & \Phi^{2}(X)
\end{tikzcd}

i.e.\ on a un foncteur co-semi-simplicial

\begin{tikzcd}[column sep=large, row sep=small, nodes={font=\scriptsize}]
\varphi \arrow[r, bend left=18] \arrow[r, bend right=18] \arrow[d, no head] &
\varphi^{2} \arrow[r, bend left=25] \arrow[r] \arrow[r, bend right=25]
  \arrow[d, no head] &
\varphi^{3} \arrow[d, no head] \\
\Phi^{0} & \Phi^{1} & \Phi^{2}
\end{tikzcd}

\textbf{NB} Les opérations \uncertain{de dégén\add{érescence}} n'utilisent que
$\beta$, \ill{} $\lambda$.

Pour tout $Y$ dans $A$, on a une co-augmentation
\[
  f(Y) \longrightarrow \Phi^{*}\bigl(f(Y)\bigr)
\]
qui est une équivalence d'homotopie, i.e.\
$f \xrightarrow{\ \approx\ } \Phi^{*} \circ f$ ; l'homotopie est donnée par
\[
  f \xrightarrow{\ f * \alpha\ } \Phi \circ f = \varphi f = fgf .
\]

Pour tout $X$ (\struck{sous-entendu} : \ill{} une donnée d'une \ill{},
\struck{\ill{}} $X \simeq f(Y)$), on considère les hom.\ possibles d'objets
co-semi-simpliciaux
\[
  X \longrightarrow \Phi^{*}(X)
\]
i.e.

\begin{tikzcd}[column sep=large, row sep=small]
X \arrow[d, "u"'] & \\
\varphi(X) \arrow[r, bend left=18, "\varphi(u)"]
           \arrow[r, bend right=18, "\lambda"'] & \varphi^{2}(X)
\end{tikzcd}

\section*{Fidélité (pages 71 à 75)}

\page{71}
Titre porté sur la chemise : « Fidélité ».

\page{72}

\begin{quote}
\textbf{Foncteurs fidèles transportables comme foncteurs oubli de structures.}
\end{quote}

\note{le premier tiers de la page est récrit deux fois et presque entièrement
biffé ; seuls les énoncés qui subsistent sont transcrits}

Soit $p : E \to B$ un foncteur. \ill{} les propriétés \add{on dit que} :

\begin{enumerate}
\item[(i)] $p$ est fibrant \struck{\ill{}} ;
\item[(i bis)] \ill{} exacte de $f^{*}$ \ill{} ;
\item[(ii)] \struck{\ill{}} pour toute flèche $x \overset{u}{\to} y$ de $B$ et
      tout $\xi \in \mathrm{Ob}\, E_{x}$, il existe $\xi \to \eta$ dont
      l'image par $p$ est $u$ ;
\item[(ii$'$)] pour toute flèche $x \to y$ de $B$, les flèches
      \struck{cocartésiennes} cartésiennes \ill{} sont des isomorphismes.
\end{enumerate}

Soit $p : E \to B$ un foncteur. On dit que $p$ est \emph{transportable} si
\struck{\ill{}} $E \times_{B} B^{is} \to B^{is}$ est fibrant (ce qui revient au
\ill{}). Donc $p$ est \struck{fibrant} transp.\ et $p^{0}$ l'est. Les flèches
cartésiennes de $E \times_{B} B^{is}$ sont les flèches cocartésiennes et
\ill{} les flèches \ill{} de $E$.

Soit $p$ un foncteur ; alors $p$ est isom.\ à un foncteur transportable.

Soient $p : E \to B$ et $p' : E \to B$ deux foncteurs transportables
isomorphes. Alors il existe \add{isomorphe à $\mathrm{id}_{E}$} une
auto-équivalence $\Phi$ de $E$ telle que $\Phi$ soit une $B$-équivalence
$(E, p) \to (E, p')$, \ill{}

\page{73}
\uncertain{conservatif} / fidèle $=$ filtre des propriétés.
\note{les deux mots sont réunis par une accolade}

\page{74}

\ill{} près.

Soit $p$ un foncteur transportable. Alors $p$ est fidèle $\Longleftrightarrow$
les catégories fibres $E_{b}$ sont des catégories ordonnées [donc si $p$ est
fibrant, $E$ est une catégorie fibrée en cat.\ ordonnées].

Pour que $p$ soit conservatif, il f.\ et s.\ que les catégories fibres soient
des \uncertain{groupoïdes} \struck{\ill{} $E \times_{B} B^{is}$ soit \ill{}}.
Donc pour que $p$ soit fidèle et conservatif, il f.\ et s.\ que les $E_{b}$
soient \emph{discrètes}.

Soit $p$ un foncteur \struck{fibrant}. \struck{Conditions équivalentes}

\begin{enumerate}
\item[a)] $p$ fibrant, fidèle, conservatif ;
\item[b)] $p$ fibrant \add{et $p$ est fidèle} à fibres discrètes ;
\item[c)] \struck{$p$ fibrant et fidèle} pour tout $X \in \mathrm{Ob}\, E$,
      $E_{/X} \to B_{/p(X)}$ est une équivalence ;
\item[d)] il existe $F \in \widehat{B}$ et une $B$-équivalence
      $E \approx B_{/F}$.
\end{enumerate}
\marginal{NB.\ la deuxième cond.\ de c) signifie que $p$ est fibrant à fibres
discrètes ?}

Les foncteurs « oubli de structure » sont les \struck{transportables} transp.\
fidèles \struck{conservatifs}. Il y a un autre fait : les catégories fibres
sont en général \struck{\ill{}} \ill{} ordonnées, i.e.\ \ill{} pré-ordonnées,
\struck{\ill{}} réduites, ce qui s'exprime en disant \struck{\ill{}} que l'on
\ill{} des relations dans la définition de la transporta\add{bilité}.

\page{75}

Si $p : E \to B$ ne satisfait pas cette condition, une \uncertain{modification
triviale} de $E$ (avec équivalence) la rend satisfaite. Partant d'un foncteur
fidèle transportable, on définit \struck{une $B^{is}$-catégorie} un foncteur
sur $B^{is}$
\[
  b \longmapsto \text{classes d'isom.\ de } E_{b}
  = \text{« structures d'espèce } p \text{ sur } b \text{ »}
\]
\add{à valeurs dans les \struck{ensembles} \ill{}} (cf.\ « \emph{transport de
structure} »). $E$ se reconstitue à isom.\ \ill{} près \ill{} objet de $B$
\ill{} muni d'une structure d'espèce $p$. De façon précise, ceci provient des
\ill{}

Le foncteur $E^{is} \overset{p^{is}}{\longrightarrow} B^{is}$ ne reconstitue
pas la catégorie $E$ elle-même. Pour pouvoir la reconstituer, il faudrait que
$E$ soit fibrant \ill{}, et on définit sur $B$, de façon précise, la notion des
« structures image inverse » d'une structure d'espèce $p$ sur un
$c \in \mathrm{Ob}\, B$, par une flèche $b \to c$, de façon à obtenir un
foncteur
\[
  \Sigma : B^{\circ} \longrightarrow (\text{ens.\ ordonnés})
\]
\add{ou $B^{is} \to (\text{ens.\ ordonnés})$}. \ill{} ainsi $E$ à équivalence
près par la connaissance de $\Sigma$, et \ill{} à isom.\ près si $E \to$ \ill{}
identité.

Ainsi, une $B$-équivalence (\uncertain{$B$-isom.}) \add{comp.\ et réduits} les
catégories fibrées sur $B$, à projection fidèle), \ill{} foncteurs
$B^{\circ} \to (\text{ens.\ ord.})$, \ill{} parlant en termes de foncteurs
\struck{\ill{}} et pour l'expression des 2-catégories, \ill{}
\struck{structures} est assez \uncertain{suggestif}.

\section*{Morphismes et comorphismes — questions sur les topos (pages 76 à 80)}

\page{76}
Titre porté sur la chemise : « Morphismes et comor\add{phismes} — Liste
Problèmes ».

\page{77}

\textbf{Questions sur les topos.} (Pour M\textsuperscript{lle} \ill{} ?)

\begin{enumerate}
\item[1)] Définition et étude des « immersions » de topos. \struck{Théorie de}
      Caractérisation ? Par la condition $f_{*}$ pleinement fidèle, i.e.\
      $f^{*} f_{*} \xrightarrow{\ \sim\ } \mathrm{id}$ (i.e.\ $f^{*}$
      \uncertain{voisin} du passage au quotient ?).
      \marginal{Théorie de $f_{!}$ et $f^{!}$}

      Dans quelle mesure les \struck{N.B.} \ill{} $f : U \to V$
      déterminent-elles la « \ill{} » ou le fermé « lui donnant naissance » ?
      Plus \ill{} et immersions fermées ? Si $U' \to V'$ \ill{} les topos
      \ill{} qui sont \ill{} : \ill{} « lieu » [i.e.\ il existe une
      équivalence $V \overset{g}{\to} V'$ rendant commutatif, à isomorphisme
      près, le diagramme des $f^{*}$, $f'^{*}$], alors $g$ est-elle déterminée
      à isom.\ unique près ? [\uncertain{Métamath.} « \ill{} de topos » ??]
      \marginal{O.K.\ pour immersions ouvertes}

      \ill{} de « monomorphisme » \ill{} Caractérisation des immersions
      ouvertes [par $f^{*}$ commutant aux limites quelconques, en plus de
      $f_{*}$ pl.\ fid.\ ?]. Caractérisation des immersions fermées [par la
      condition $F \to f_{*} f^{*} F$ un épi pour \uncertain{$H$} faisceau
      abélien \ill{} — \ill{}, autre condition, en plus de $f_{*}$ pl.\ fid.]?
      \marginal{Application aux \ill{} $F = \emptyset$ \ill{} « \ill{} $H^{i}$
      aux $\varinjlim$ »}
      \ill{} $f^{*}$ commute aux \ill{} quelconques.

\item[2)] Étude des $f$ tels que $f$ \ill{}. Alors $f^{*}$ \struck{\ill{}}
      admet-il un adjoint à gauche $f_{!}$ ? Dans quelle mesure le \ill{}
      $f^{*} : U \to U_{/S}$ \ill{} ces deux foncteurs \ill{} est-il typique ??
      \marginal{Il n'est ni l'un ni l'autre, cf.\ p\textsuperscript{te}
      essentiel.\ Regarder \ill{} la commutation aux \emph{limites} plutôt
      \ldots}
\end{enumerate}

\page{78}

\begin{enumerate}
\item[3)] Étude du cas où $f^{*}$ pl.\ fidèle. Étude du cas où $f^{*}$ commute
      aux $\varprojlim$ \emph{finies}. [et aussi, à \ill{} on transfère
      \ill{}]

\item[4)] Étude du cas où $f_{*}$ pl.\ fidèle (cf.\ 1°)) \ill{}

\item[5)] Notion de morphisme \emph{dominant} de topos
      ($f_{*}(\emptyset) = \emptyset$). \struck{Plus} Lien avec le fait
      $f^{*}$ \emph{conservatif}. \ill{} $f_{*}$ conservatif ? (Peut-on dire
      \ill{} chose sur le cas $f_{*}$ conservatif ?)

\item[6)] Donner une caractérisation des foncteurs exacts $U \to V$ provenant
      \ill{} i.e.\ \ill{} Cartier-Artin. (Cf.\ Deligne.)

\item[7)] Les \struck{morphismes} de $\underline{U}$-topos forment une
      $U$-cat\add{égorie} \ill{} (i.e.\ les $\mathrm{Hom}$ sont « petits »).
      L'ens.\ des classes, à isom.\ près, de tels morphismes est-il petit ?
      \marginal{non}

\item[8)] a) Théorie de \ill{} :
      \marginal{topos noethériens ?}

      \begin{tikzcd}[column sep=huge]
      (\mathrm{Top}) \arrow[r] &
      (\mathrm{Cat}) \arrow[l, bend left=30, "C \mapsto \widehat{C}"]
      \end{tikzcd}

      b) Topos \ill{} $H^{i}$ \ill{} \ill{} finies \ill{}
      \marginal{on y a\ill{} lieu de considérer \ill{} une catégorie \ill{}}

\item[9)] Définir \struck{\ill{}} les \emph{catégories} \emph{fibrées} de
      topos \ill{} dans certains cas, la situation. (Cas dont
      \uncertain{obtenir} de nouvelles catégories \ill{} discrètes).
      \marginal{ne serait-ce pas plus explicite que (Cat) ?}

\item[10)] Topos \ill{}. Catégories fibrées resp.\ champs \ill{}
      essentiellement représentables \ldots

\item[11)] \begin{tikzcd}[column sep=huge]
      (\mathrm{Ens.}) \arrow[r] &
      (\mathrm{Esp.}) \arrow[r] &
      (\mathrm{Top}) \arrow[ll, bend left=30]
      \end{tikzcd}
      \note{le qualificatif qui suit « Esp.\ » n'a pas pu être lu}

      \marginal{2-adjoint à \ill{} des foncteurs \emph{fibres}}
\end{enumerate}

\page{79}\note{page barrée de longs traits diagonaux}

\ill{} deux adjoints de $f$ coïncident et sont isom.

Il suffirait de l'avoir pour \struck{deux} catégories \struck{\ill{}} $C$,
$C'$, telles que $f$, $g$ \ill{} pl.\ fid.\ \ill{} soient \ill{} par des
\struck{foncteurs} \ill{} : $\widehat{C}$, $\widehat{C'}$). Peut-on \ill{}
prendre $C = C'$, $f = g$, \ill{} ? Peut-on \ill{} donner une catégorie $C$, et
\ill{} $f : C \to C$ qui \ill{} soit \ill{} par une \ill{} pl.\ fid.\ \ill{},
et qui soient isom.\ \ill{} adjoint (à droite et à gauche)) ?
\marginal{O.K.}

Prendre $C$ une catégorie ayant un objet $e$, mais telle que le foncteur
$f : C \to C$ défini par $f(x) = e$ \ill{} constant \ill{} $C \to C$ ne soit
pas pl.\ fidèle, i.e.\ $C$ \ill{} une catégorie discrète, dans $C$ ayant deux
\ill{} $x \neq e$.

[Ex.\ minimal :

\begin{tikzcd}[column sep=huge]
e \arrow[r, bend left=25, "u"] & a \arrow[l, bend left=25, "v"]
\end{tikzcd}

et $w : a \to a$ : deux objets $e$, $a$, trois flèches $u$, $v$, $vu = w$,
\ill{} flèches identiques \ill{} composition \ill{} déterminées uniquement
\ill{} les relations $[uv = \mathrm{id}_{e}$, $uw = u$, $wv = v]$, sauf
\ill{} $w^{2} = \mathrm{id}_{a}$ \ill{} : pour \ill{} il faut dire ce qu'est
$w^{2}$ : on peut prendre $w^{2} = w$ \ill{} ou deux]. On trouve le \ill{}

\page{80}

\ill{} adjoint à droite \ill{}, \ill{} en \ill{} les \ill{} itérées

\begin{tikzcd}[column sep=small, row sep=small, nodes={font=\scriptsize}]
f_{!} & f^{*} \arrow[d, no head] & f_{*} \arrow[d, no head] &
  f^{!} \arrow[d, no head] & f_{?} \arrow[d, no head] \\
& g_{!} & g^{*} \arrow[d, no head] & g_{*} \arrow[d, no head] &
  g^{!} \arrow[d, no head] \\
& & h_{!} & h^{*} & h_{*}
\end{tikzcd}

etc. On trouve ainsi les résultats \add{de} \ill{}

\textbf{Proposition.} Soit $U$ un univers, \struck{\ill{}} contenant un
ensemble infini. Pour toute catégorie \ill{} il existe deux $U$-topos $E$, $F$
[de la forme $\widehat{C}$, $\widehat{C'}$] \ill{} avec \ill{} ($C$, $C'$ des
\emph{petites} catégories) et une \ill{} \add{dont aucun n'est pl.\ fid.})
\ill{} des $n$ foncteurs ($f_{1}, f_{2}, f_{3}, \ldots, f_{n}$ \ill{})
alternativement de $E$ dans $F$ et de $F$ dans $E$, chacun étant l'adjoint à
droite du \ill{} (\ill{} vaut \ill{} les $f_{i}$) qui précède \ill{}
$g$ \ill{} $= f_{1} \cdots f_{n}$, et \ill{} un \uncertain{morphisme} \ill{}
\ill{} un \uncertain{morphisme direct}].

\textbf{Question.} Peut-on trouver une suite infinie dans les deux
directions ? Il suffirait de trouver \ill{} $f$ et $g$ tels que $g$ soit à la
fois adjoint \ill{} \ill{} à droite de $f$, i.e.\ tels que les
\note{la page s'arrête au milieu de la phrase ; la suite est au lot 5}

\end{document}
